\section{Исходный код}
Программа \texttt{main.py} состоит из класса \texttt{Crawler} и вспомогательных
функций. Основные части: \texttt{normalize} (нормализация URL),
\texttt{wanted} (фильтр ссылок), \texttt{enqueue}, \texttt{requeue}, \texttt{next\_item}
(очередь), \texttt{crawl\_one} (загрузка и сохранение одной страницы),
\texttt{finish} (учёт результата) и \texttt{run} (основной цикл).

Нормализация URL (\texttt{SAFE} --- символы, которые MediaWiki не кодирует в заголовках) и фильтр ссылок (домен, путь, пространство имён, robots.txt):

\begin{lstlisting}[language=Python]
def normalize(url, base=None):
    p = urlparse(urljoin(base, url) if base else url)
    path = unquote(p.path).replace(' ', '_')
    return 'https://' + p.netloc.lower() + quote(path, safe=SAFE)

def wanted(self, url):
    p = urlparse(url)
    if p.netloc != self.domain or not p.path.startswith('/wiki/'):
        return False
    title = unquote(p.path[len('/wiki/'):])
    ns, colon, rest = title.partition(':')
    return (not (colon and ns in self.skip)
            and self.robots.can_fetch(self.ua, url))
\end{lstlisting}

Очередь в MongoDB: добавление без повторов, выбор следующего URL, возврат в
очередь устаревших документов:

\begin{lstlisting}[language=Python]
def enqueue(self, urls):
    now = time.time()
    ops = [UpdateOne({"_id": u}, {"$setOnInsert": {"state": "new",
           "tries": 0, "refresh": False, "added": now}}, upsert=True)
           for u in set(urls) if self.wanted(u)]
    if ops:
        self.queue.bulk_write(ops, ordered=False)

def next_item(self):
    flt = {"state": "new"}
    if self.limit and self.docs.estimated_document_count() >= self.limit:
        flt["refresh"] = True      # corpus is full: only refresh
    return self.queue.find_one(
        flt, sort=[("tries", ASCENDING), ("added", ASCENDING)])

def requeue(self):
    cutoff = int(time.time() - self.refresh_days * 86400)
    self.queue.update_many({"state": "done", "done_at": {"$lt": cutoff}},
        {"$set": {"state": "new", "tries": 0, "refresh": True}})
    self.queue.update_many({"state": "error"},
        {"$set": {"state": "new", "tries": 0}})
\end{lstlisting}

Загрузка одной страницы: условный запрос, разбор ошибок HTTP, хэш текста
статьи, запись документа с историей версий:

\begin{lstlisting}[language=Python]
def crawl_one(self, url):
    old = self.docs.find_one({"url": url}, {"html_content": 0})
    headers = ({"If-Modified-Since": old["last_modified"]}
               if old and old.get("last_modified") else {})
    try:
        r = self.session.get(url, headers=headers, timeout=self.timeout)
    except requests.RequestException:
        return "error"
    if r.status_code == 304:                    # not modified
        self.docs.update_one({"url": url},
            {"$set": {"crawl_timestamp": int(time.time())}})
        return "done"
    if r.status_code == 429 or r.status_code >= 500:
        return "error"                          # retry later
    if r.status_code != 200 or "text/html" not in r.headers.get(
            "Content-Type", ""):
        return "failed"                         # never retry
    soup = BeautifulSoup(r.text, "html.parser")
    body = soup.select_one("#mw-content-text") or soup
    text = body.get_text(" ", strip=True)
    h = hashlib.md5(text.encode()).hexdigest()
    ...                                         # canonical URL, see src
    if old and old["content_hash"] == h:        # same text: only the date
        self.docs.update_one({"url": doc_url},
            {"$set": {"crawl_timestamp": now}})
    else:
        update = {"$set": {"url": doc_url, "html_content": r.text,
                  "source": self.name, "crawl_timestamp": now,
                  "content_hash": h, ...}}
        if old:                                 # changed: keep history
            update["$push"] = {"history": {
                "content_hash": old["content_hash"],
                "crawl_timestamp": old["crawl_timestamp"]}}
        self.docs.update_one({"url": doc_url}, update, upsert=True)
    self.enqueue(normalize(a["href"], url) for a in body.select("a[href]"))
    return "done"
\end{lstlisting}

Основной цикл и остановка по сигналу (очередь в базе --- контрольная точка):

\begin{lstlisting}[language=Python]
def run(self):
    self.requeue()
    self.enqueue(normalize(s) for s in self.seeds)
    while not self.stop.is_set():
        item = self.next_item()
        if not item:
            break
        self.finish(item, self.crawl_one(item["_id"]))
        self.stop.wait(self.delay)    # returns at once on Ctrl+C
\end{lstlisting}

\pagebreak
\section{Выводы}
Выполнив вторую лабораторную работу по курсу \enquote{Информационный поиск}, я
написал поискового робота для русских Викиновостей. Основное, что я узнал
нового:

\begin{itemize}
    \item устройство краулера на практике: очередь URL, нормализация адресов,
    обход в ширину, вежливые задержки, \texttt{robots.txt}, идентифицирующий
    \texttt{User-Agent};
    \item очередь удобно хранить в самой базе данных: тогда она одновременно
    защищает от повторных загрузок (\texttt{upsert} по URL) и служит
    контрольной точкой --- остановить робота можно в любой момент, не заводя
    отдельного файла состояния;
    \item ошибки HTTP нужно делить на постоянные (404, не-HTML) и временные
    (429, 5xx, таймауты): первые не повторяются, вторые откладываются в конец
    очереди с ограничением числа попыток, и ни одна не должна останавливать
    всю обкачку;
    \item переобкачка по изменению: условный запрос \texttt{If-Modified-Since}
    экономит трафик, а сравнение хэша текста статьи позволяет не переписывать
    неизменившиеся документы и хранить историю версий изменившихся;
    \item реальность внесла поправки: заголовки новостей содержат двоеточие,
    поэтому служебные страницы (\texttt{Категория:}, \texttt{Служебная:}, \dots)
    нельзя отсеивать по одному символу --- понадобился список пространств
    имён; одна и та же страница доступна по нескольким адресам, поэтому
    документы хранятся под каноническим URL (в контрольном прогоне такими
    оказались 22 из 222 адресов);
    \item «сырой» HTML страницы Викиновостей в 17 раз больше полезного текста
    (125{,}7 КБ против 7{,}5 КБ на документ), что нужно учитывать при выборе
    хранилища и при построении индекса.
\end{itemize}

Викиновости открыты для автоматического доступа при соблюдении
правил (идентифицирующий \texttt{User-Agent}, \texttt{robots.txt}, пауза между
запросами), поэтому для робота выбран именно этот источник. Скорость обкачки
упирается в паузу между запросами (1~с): корпус из десятков тысяч документов
набирается за часы. Собранные документы с датами обкачки и историей версий
будут использоваться в следующих лабораторных работах для построения индекса.

\pagebreak
